\section{Anthropic、Neo Labs 与模型公司分化}

前一章把 OpenAI 拆成产品和现金流，本章看模型公司之间的分化。Anthropic 与 OpenAI 的差异，不只是模型名字不同，而是客户结构、产品路线、毛利率控制和财务预测假设不同。Neo Labs 则代表另一种分化：从大模型公司溢出的研究团队，继续押 Continuous Learning、定制训练和开源生态。

\begin{figure}[H]
\centering
\includegraphics[width=0.94\textwidth]{openai-vs-anthropic.png}
\caption{OpenAI vs Anthropic：2C 产品公司与 2B/Coding 公司。自制概念图，依据 00:28:25--00:31:25 对谈内容整理。}
\end{figure}

\begin{knowledgebox}{读图：两家公司不是同一种生意的复制}
OpenAI 的优势更偏 C 端心智、个人助理和企业入口；Anthropic 超过多数收入来自 2B，更强调 Coding、Finance、Agent 模型和企业知识工作者。二者都可以很好，但收入质量、销售路径和产品组织方式不同。
\end{knowledgebox}

\subsection{Anthropic 的增长与毛利率控制}

Freda 提到，Anthropic 的收入从 10 亿到 70 亿年化增长很快，最近几个月和 OpenAI 每月新增年化收入在数量上接近。更重要的是毛利率改善：通过模型架构、auto scaling、推理优化等方式，模型公司毛利率并非注定很低。她认为长期模型公司毛利率到 70\%--80\% 并非不可想象。

这里要注意财务预测的假设差异。她说 Anthropic 的预测似乎假设 2028 年后总算力不再大幅增长，通过训练成本增速下降拉升利润率；OpenAI 的预测则更像假设模型 ROI 逐年提升，即投入一块钱训练，未来能收回更多倍数。这两个假设都不是事实本身，而是值得投资人追问的模型。

\begin{importantbox}{模型公司财务预测的两个按钮}
第一是成本按钮：训练和推理成本增速是否下降。第二是收入按钮：每一代模型能否带来更高 ROI。不同公司可以通过不同按钮讲出利润率改善故事。
\end{importantbox}

\subsection{Neo Labs：下一批模型公司的融资逻辑}

本节把视角从 OpenAI/Anthropic 两强扩展到新实验室。前面说模型公司会分化，Neo Labs 的问题就是：在头部模型公司已经很强时，新团队还能靠什么融资、靠什么差异化？

节目里提到 SSI、Thinking Machines 等熟悉名字，也提到最近几个月又有多家 Neo Labs 融资，很多由学术背景人员创立，融资金额很大。Freda 的疑问是：大模型是否最终会收敛到头部一两家，剩余公司会转向拥抱开源？Thinking Machines 第一个产品 Thinker 被描述为帮助用户定制训练开源模型，这背后其实是在押开源市场越来越大。

\begin{knowledgebox}{术语消化：Neo Labs 与开源定制}
\begin{tabularx}{\textwidth}{p{0.24\textwidth}p{0.32\textwidth}X}
\toprule
术语 & 一句话解释 & 本期关系 \\
\midrule
Neo Labs & 新一代模型/研究实验室，常由大厂或学术人才创立 & 承接下一范式、开源定制和专业模型机会。 \\
Continuous Learning & 持续学习，从新交互和数据中改进模型 & SSI 被节目提到可能在探索这一方向。 \\
开源定制训练 & 在开源模型基础上为用户做专门训练 & Thinking Machines 的 Thinker 被用来说明这一方向。 \\
OpenAI 对标估值 & 用 OpenAI 的总估值作为其他模型公司相对锚点 & 一级市场模型公司融资常用的心理标尺。 \\
\bottomrule
\end{tabularx}
\end{knowledgebox}

\subsection{本章小结}

Anthropic 和 Neo Labs 说明模型公司已经分化：有的做 C 端入口，有的做企业/Coding，有的做开源定制，有的押下一范式。资本市场会不断用 OpenAI 做锚，但真正的差异来自客户结构、毛利率控制和产品路线。

